\documentclass[11pt]{article}
\usepackage[margin=1in]{geometry}
\usepackage{amsmath,amssymb}
\setlength{\parindent}{0pt}
\setlength{\parskip}{6pt}
\begin{document}
\section*{STAT 412 QZ08 --- Question 5}

\subsection*{Solution}

Let \(\mu_1\) be the mean tip percentage with no candy and \(\mu_2\) be the mean tip percentage with two candies.

\[
H_0:\mu_1=\mu_2,\qquad H_1:\mu_1<\mu_2.
\]

Using Welch's test,
\[
\text{SE}=\sqrt{\frac{1.46^2}{25}+\frac{2.34^2}{25}}=0.5518,
\]
\[
t=\frac{18.14-20.58}{0.5518}=-4.42.
\]
The Welch degrees of freedom are
\[
df\approx40.23.
\]
\[
p\approx0.000.
\]

Since \(p<0.01\), reject \(H_0\). There is sufficient evidence that giving candy results in greater tips.

A confidence interval suitable for the one-tailed \(\alpha=0.01\) test is the 98\% two-sided interval:
\[
(\bar{x}_1-\bar{x}_2)\pm t^*SE.
\]
Using \(t^*\approx2.423\) with \(df\approx40.23\),
\[
-2.44\pm 2.423(0.5518)=(-3.78,-1.10).
\]
\[
-3.78<\mu_1-\mu_2<-1.10.
\]
Because the entire interval is less than 0, it supports the hypothesis-test conclusion.

\subsection*{Final Answer}

\fbox{\begin{minipage}{0.94\linewidth}
\(t=-4.42\), \(p=0.000\). Reject \(H_0\). CI blanks: lower \(-3.78\), upper \(-1.10\). The 98\% CI is \((-3.78,-1.10)\), which supports the conclusion.
\end{minipage}}

\end{document}
